\section{Discussion}
\label{sec:discussion}
\paragraph{What the evidence says} The gains that learning-guided search reports are mostly \emph{structure} (destruction--reconstruction), \emph{budget} and \emph{tuning}. The learner's own contribution, the decision effect, is bounded by one quantity, the headroom of the best fixed action net of run-to-run noise: below about $0.05$ points the online learner never gained beyond the equivalence margin in any cell outside those that set the threshold, and above it none was shown to once the order of the actions is randomised (the standard learner repeats the action listed first, and an order-invariant Q-Learning gains beyond the margin in none of $85$ cells). What headroom rewards is knowledge acquired \emph{before} the run: a validated action, a selector trained offline (about four fifths of the one-seed oracle gain on $P||\Cmax$, $60$--$73\%$ on unrelated machines) or a neural policy (two thirds of the validated gain; $60\%$ out of that $67\%$ is already a constant mixture). How much offline knowledge is needed depends on the host: on the standard $R||\Cmax$ benchmark one action for all instances already gains $51\%$ of the class-validated gain and the size of the instance supplies $89\%$ of it, whereas in the TSP no finer knowledge adds anything. The one regime in which an online learner beat an offline static action is a transferred action applied to an unseen extreme size; there the learner is an insurance, best built as a correction of the offline prior: warm-started with it, the optimistic Q-Learning gains $0.67$ points over the same prior without learning when the prior was transferred and loses $0.28$ when it was exact. Where actions are complementary (the genetic algorithm, the ILS operators) the static competitor must be a \emph{mixture}; the claim that a validated static configuration is never beaten beyond the margin concerns the best of a single action or a validated mixture.

\paragraph{Why the usual comparison misleads} Comparing a learning-guided metaheuristic with a non-learning one changes at least four things at once: the search structure, the number of iterations or the CPU time, the tuning effort of each side, and the information the learner receives. The protocol varies one at a time: the learning-guided IG appeared to owe its advantage over the ILS to the learner when it belongs to the destruction--reconstruction mechanism; at equal CPU time the learner appeared responsible for a gain that belongs to choosing cheaper actions; and a one-seed oracle appeared to show room for improvement that was mostly noise. In the literature we audited, $22$ of $23$ papers report a benefit of the learned component, $10$ include a random-selection control, $6$ compare at a common time limit and $3$ a static configuration; the audit counts which controls were reported, not whether the claims are true.

\paragraph{A checklist for a claim that a learned selector helps}
\begin{enumerate}
\item Report a \textbf{random control over the same actions at the same budget}, with the same code path as the learner.
\item At equal CPU time, \textbf{calibrate the cost ratio} $K$ on separate instances, report it for every arm, and split the gain into a decision effect and a cost effect.
\item \textbf{Validate a static action (or a mixture) offline} on instances disjoint from the test set, with the same effort the learner's tuning received; add a configurator and a published AOS design where they apply.
\item Estimate the headroom \textbf{with cross-fitting}, and report the one-seed value (an upper bound) and the cross-fitted value (a lower bound).
\item Include a \textbf{positive control} and an \textbf{equivalence margin}, so that ``no effect'' is a finding and not a failure to reject.
\item \textbf{Record hypotheses, decision rules and predictions before running}, with hashes of the code and frozen inputs, and control the error rate within each family of tests.
\item Report the learner's \textbf{sensitivity to the order of the actions}: the share of choices on the action listed first, and the result with the actions permuted (or ties broken at random).
\item State \textbf{what the static competitor knows} and test its transfer to an instance type or size it was not tuned for; an online learner warm-started with the offline prior can be worth its cost exactly where the transferred action fails, and costs something where the prior is exact.
\end{enumerate}

\paragraph{What the registered predictions showed} \npredheld{} of \npredtotal{} point predictions held in \nbatteries{} batteries. The misses (cost ratios, the classic AOS methods, the flowshop, warm-started Q-Learning, finer times) led to the headroom rule and to cross-fitting. The headline pairs keep their significance under Pratt's treatment of ties and a permutation test, $1{,}606$ of the $1{,}674$ confirmed tests survive a global Benjamini--Yekutieli control over the $3{,}578$ tests of all registered batteries (each extension of the control was itself registered, three predictions of three).

\section{Limitations}
\label{sec:limitations}
\begin{enumerate}[label=(\roman*)]
\item \emph{Equal CPU time} depends on $K$, measured on separate instances in sequence and reported for every arm; a learner whose actions are cheap can win on cost alone, which is why the decomposition is part of the protocol.
\item \emph{Practical significance.} Gains are in points of the makespan, and the margin is a choice. The hosts are not far from the best known values: on the standard benchmark the validated action ends $4.89\%$ above the best known solutions at $T=1000$ and $2.69\%$ at $T=10^4$, and in the TSP $1.21\%$ above the optimum of the TSPLIB instances (T15).
\item \emph{Registration and multiplicity.} Each battery has its own Holm family; the headroom rule began as an exploratory reading. The pre-specification is a dated, hash-stamped diary under version control, not a public registry.
\item \emph{Scope.} The evidence is for $P||\Cmax$, $R||\Cmax$, a permutation flowshop, one genetic algorithm and one TSP local search, with tabular and bandit learners, a DQN on joint actions and two published AOS designs. Deep policy gradients that learn during the run and portfolios of hundreds of operators are outside the evidence, and the neural policy is one architecture trained on one host.
\item \emph{Implementation dependence.} A conclusion about ``Q-Learning'' is a conclusion about an implementation (initial values, tie rule, reward sign); we did not rerun UCB1, the policy gradient and the roulette with permuted actions, and the irace-tuned Q-Learning keeps the standard tie rule.
\item \emph{Offline knowledge.} The map by granularity ran on one scheduling benchmark and the TSP; the unseen-size reading is post hoc and rests on the two extreme sizes of four; the warm-start result is for one learner and one prior construction.
\item \emph{Horizons} up to $T=10^5$ were studied on two hosts and $T=10^4$ on the standard benchmark in the richest space only; a learner that needs a longer horizon than the headroom lasts for is not excluded in general.
\item \emph{Selector cost.} Training costs of the offline selector and of the neural policy ($115{,}200$ and $138{,}240$ IG runs, eight to fourteen times the tuning of the validated action) are not charged to them.
\item \emph{Review status.} The manuscript has not been externally reviewed; simulated blind reviews guided the revisions but are not a substitute.
\end{enumerate}

\section{Conclusion}
\label{sec:conclusion}
A learned operator selector is not what explains the advantage of an IG over ILS for identical parallel machines: that belongs to destruction--reconstruction, and the share of a local search's gain that is only budget ($93\%$, $30\%$ or $28\%$ by implementation) makes the CPU cost ratio a result to report. The oracle headroom net of run-to-run noise bounds what a validated action or an offline policy can capture, and online Q-Learning did not pay at these budgets: the gains first attributed to it came from the standard implementation repeating the action listed first, and an order-invariant version gains beyond the margin in none of $85$ cells. A validated static configuration was never worse than an online learner beyond the margin on any host, on published instance families, on the standard benchmark and at a ten times longer horizon (at equal CPU time on the standard benchmark it lost two of $49$ pairs). What pays is offline knowledge: one action chosen offline for all instances beats every online learner, the size of the instance carries nine tenths of what class-level knowledge gives, and an online learner is worth having as a corrector of that knowledge where it does not cover the instance (a size never tuned for): warm-started with the offline prior it recovers $48\%$ of the class-level gain there, against $17\%$ for the transferred static action, and without a prior it recovers $14\%$. The same asymmetry appears in the TSP, an iterated local search for a different kind of problem: learning helps against a wrong prior and costs little against a nearly exact one.

For the practitioner the message is a procedure more than a verdict: before crediting a learned component, compare it with random choice at the same budget, with a static choice validated offline and with a configurator, permute the order of the actions to see whether the learner is learning or repeating the first item of its list, and ask whether the headroom is large enough for any online method to find. For the researcher it is a set of estimators with stated and checked properties and a registry of $\npredtotal$ pre-recorded predictions, whose misses are listed and changed the analysis.

\paragraph{Data and code availability} Code, instances, results, the research diary with the pre-registered hypotheses and md5 hashes, the protocol library and the scripts that recompute every number are at \url{https://github.com/CarlosAraujo7/pmsp}. \todonote{funding statement: CAPES/FUNCAP/CNPq to confirm; anonymised repository link for review}
